\section{Persistenza dei Dati e Database Relazionali}
I sistemi di basi di dati (DBMS) rappresentano l'infrastruttura standard per l'archiviazione strutturata, la persistenza e il recupero delle informazioni nei moderni sistemi \textit{software}.

Il paradigma principale per l'organizzazione logica dei dati è il modello relazionale~\autocite[ch. 2]{atzeniBasiDiDati2009}, introdotto formalmente da Edgar F. Codd nel 1970. Questo modello, a differenza dei modelli gerarchici o reticolari precedenti, si basa sulla teoria matematica degli insiemi e sull'algebra relazionale. In questo modello, le informazioni sono organizzate in tabelle, chiamate relazioni, composte da righe, o tuple, che rappresentano singole istanze di un'entità, e colonne, o attributi, che descrivono le proprietà di quell'entità. La solidità matematica del modello consente di interrogare e manipolare i dati in modo dichiarativo, separando la logica applicativa dalle modalità fisiche di memorizzazione delle informazioni. Tale separazione permette una maggiore indipendenza tra i processi di sviluppo \textit{software} e le tecniche di gestione dello spazio di archiviazione.

\subsection{Sistemi di Gestione di Database}
La gestione del modello relazionale è affidata ai sistemi di gestione di basi di dati (DBMS). I DBMS sono sistemi \textit{software} avanzati progettati per gestire collezioni di dati, garantendo che tali dati siano condivisi, persistenti e mantenuti con elevati standard di affidabilità e \textit{privacy}. Queste collezioni di dati possono essere di dimensioni molto grandi, rendendo necessaria la gestione dell'allocazione dei dati in memorie secondarie. Il DBMS garantisce la condivisione tra molteplici utenti e applicazioni, eliminando la ridondanza strutturale e prevenendo il rischio di inconsistenze logiche derivanti dalla replicazione non controllata. Questo viene implementato tramite rigorosi meccanismi di controllo della concorrenza da parte dei DBMS, i quali governano gli accessi simultanei garantendo stabilità transazionale anche su volumi di milioni di operazioni al secondo. Infine, il DBMS garantisce l'affidabilità dei dati, che possono essere ripristinati pure in caso di guasti \textit{hardware} o \textit{software} (\textit{backup} e \textit{recovery}), e di sicurezza degli accessi tramite meccanismi di autorizzazione che abilitano gli utenti solo a specifiche azioni~\autocite[ch. 1.2]{atzeniBasiDiDati2009}.

L'architettura interna di un DBMS relazionale (RDBMS) si compone di vari sottosistemi: il \textit{Query Processor}, che analizza, ottimizza ed esegue le interrogazioni formulate nel linguaggio standard SQL (\textit{Structured Query Language}), lo \textit{Storage Manager}, che interfaccia il motore logico con il \textit{file system }del sistema operativo per l'allocazione dello spazio su disco e il \textit{Transaction Manager}, responsabile del controllo della concorrenza e dell'affidabilità dei dati. La sicurezza e il controllo degli accessi sono invece demandati ai moduli di autorizzazione del sistema. L'impiego di un DBMS solleva lo sviluppatore dall'onere di gestire l'accesso a basso livello ai file, la frammentazione dei dischi e le collisioni derivanti dall'accesso simultaneo ai dati da parte di molteplici processi~\autocite[ch. 1]{silberschatzDatabaseSystemConcepts2020}.

\subsection{Eloquent ORM}
L'\textit{Object-Relational Mapping} (ORM)~ è un paradigma architetturale ideato per colmare l'incongruenza strutturale e logica che esiste fra la programmazione orientata agli oggetti, come PHP, e i RDBMS. Eloquent~\autocite{Eloquent} è la soluzione ORM integrata in Laravel, sviluppata secondo il \textit{design pattern} \textit{Active Record}. In conformità a questo \textit{pattern}, ogni tabella del database viene associata a una classe \textit{Model}, le cui istanze incapsulano simultaneamente i dati della singola tupla e la logica necessaria per l'accesso e la manipolazione del \textit{record} corrispondente.

L'architettura di Eloquent si basa sul principio di convenzione rispetto alla configurazione: se lo sviluppatore non specifica esplicitamente il nome della tabella di riferimento, il sistema lo deduce automaticamente convertendo il nome della classe del modello dal formato singolare al plurale in \textit{snake\_case}. Allo stesso modo, Eloquent assume sempre la presenza di una chiave primaria intera e autoincrementante chiamata \texttt{id} e di colonne \textit{timestamp}, \texttt{created\_at} e \texttt{updated\_at}, che l'ORM gestisce automaticamente durante le operazioni di inserimento e aggiornamento.

Inoltre, Eloquent astrae le interazioni tra le tabelle del database implementandole come metodi all'interno delle classi modello. Secondo la documentazione tecnica, il sistema gestisce automaticamente le relazioni standard, come quelle uno-a-uno (\textit{hasOne}) e uno-a-molti (\textit{hasMany}), oltre alle associazioni molti-a-molti, manipolando le tabelle intermedie senza richiedere l'intervento diretto dello sviluppatore~\autocite{EloquentRelationshipsLaravel}.

\subsection{Versionamento del Database: Migrations}
Accanto al sistema ORM Eloquent, Laravel mette a disposizione un meccanismo denominato \textit{Migrations}~\autocite{DatabaseMigrationsLaravel}, progettato per gestire in modo strutturato e versionato l'evoluzione dello schema del database. Le \textit{migrations} possono essere considerate come un sistema di controllo delle versioni per la struttura della base di dati, analogo ai sistemi di versionamento del codice sorgente, ma applicato alle tabelle, ai vincoli e agli attributi del database.

Tradizionalmente, le modifiche allo schema vengono effettuate mediante l'esecuzione diretta di istruzioni SQL sul database. Tale approccio può risultare problematico nei progetti collaborativi o in sistemi soggetti a frequenti evoluzioni, poiché rende complessa la sincronizzazione delle modifiche tra ambienti differenti e aumenta il rischio di inconsistenze strutturali. Le \textit{migrations} affrontano questo problema rappresentando ogni modifica come una classe PHP che descrive esplicitamente le operazioni da applicare.

Ogni \textit{migration} implementa generalmente due metodi principali: \texttt{up()} e \texttt{down()}. Il metodo \texttt{up()} definisce le modifiche da applicare allo schema, come la creazione di nuove tabelle, l'aggiunta di colonne o la definizione di relazioni; il metodo \texttt{down()} descrive invece le operazioni inverse necessarie per annullare tali modifiche. Questo approccio consente di eseguire procedure di \textit{rollback}, riportando il database a uno stato precedente nel caso di errori o necessità di ripristino.

Laravel mette inoltre a disposizione uno strato di astrazione chiamato \textit{Schema Builder}, che consente di definire le strutture del database attraverso un'interfaccia orientata agli oggetti, riducendo la dipendenza da specifiche implementazioni SQL. In questo modo, lo sviluppatore può descrivere lo schema in maniera dichiarativa, mantenendo una maggiore portabilità tra differenti sistemi di gestione di database supportati dal \textit{framework}.

A completamento del versionamento strutturale, Laravel integra il meccanismo dei \textit{Seeders}~\autocite{DatabaseSeedingLaravel}. Mentre le \textit{migrations} operano a livello di definizione dello schema, i \textit{seeders} gestiscono il popolamento automatizzato delle tabelle mediante \textit{dataset} iniziali o di esempio. Tale meccanismo separa logicamente la struttura del database dai dati in esso contenuti, permettendo di automatizzare la configurazione dell'ambiente applicativo e garantendo la riproducibilità dello stato iniziale del sistema durante le fasi di sviluppo, \textit{testing} e distribuzione. L'impiego delle \textit{migrations} favorisce quindi la tracciabilità delle modifiche, la riproducibilità degli ambienti di sviluppo e produzione e una gestione più controllata dell'evoluzione del database durante l'intero ciclo di vita del \textit{software}.
